% preamble.tex -- shared preamble of ulm-lam-companion.tex (LuaLaTeX)

\usepackage[a4paper,inner=2.6cm,outer=2.4cm,top=2.6cm,bottom=2.8cm,headheight=14pt]{geometry}

% ---------------------------------------------------------------- fonts
\usepackage{amsmath}
\usepackage{fontspec}
\usepackage{unicode-math}
\setmainfont{Libertinus Serif}
\setsansfont{Libertinus Sans}
\setmathfont{Libertinus Math}
\setmonofont{JuliaMono}[Scale=MatchLowercase]

% ------------------------------------------- paragraphs: no indent, \medskip
\usepackage[skip=\medskipamount]{parskip}

% ---------------------------------------------------------------- packages
\usepackage[dvipsnames]{xcolor}
\usepackage{microtype}
\usepackage{booktabs,tabularx,array,xltabular}
\usepackage{enumitem}
\usepackage{needspace}
\usepackage{listings}
\usepackage[most]{tcolorbox}
\usepackage{tikz}
\usetikzlibrary{positioning,arrows.meta,fit,calc,backgrounds,decorations.pathreplacing}
\usepackage{titlesec}
\usepackage{fancyhdr}
\usepackage[hyphens]{url}
\usepackage[unicode,bookmarksnumbered,linktoc=all]{hyperref}
\usepackage[nameinlink,capitalise,noabbrev]{cleveref}

\definecolor{accent}{HTML}{1F5A96}
\definecolor{metacolor}{HTML}{555555}
\definecolor{bg}{gray}{0.95}      % code blocks
\definecolor{ttbg}{gray}{0.92}    % inline typewriter text
\definecolor{theorybg}{HTML}{EAF2FB}
\definecolor{limitbg}{HTML}{FBF0E6}
\definecolor{kw}{HTML}{1F4E9A}
\definecolor{cm}{HTML}{6B6B6B}
\definecolor{ty}{HTML}{2E6B30}
\definecolor{st}{HTML}{8A3B12}
\hypersetup{colorlinks=true,linkcolor=accent,urlcolor=accent!70!black,citecolor=accent,
  pdftitle={ulm.lam and the Field Guide: A Companion},pdfauthor={}}

\setlist{itemsep=2pt,topsep=3pt,parsep=0pt}
\renewcommand{\arraystretch}{1.18}
\emergencystretch=3em
\raggedbottom

% ---------------------------------------------------------------- headings
\titleformat{\chapter}[display]
  {\sffamily\bfseries\huge\raggedright}
  {\large\color{accent}\MakeUppercase{\chaptertitlename}~\thechapter}{0.4em}{}
\titlespacing*{\chapter}{0pt}{0pt}{1.8em}
\titleformat{\section}
  {\sffamily\bfseries\Large\raggedright}{\color{accent}\thesection}{0.8em}{}
\titlespacing*{\section}{0pt}{2.0em plus 0.5em minus 0.3em}{0.5em}
\titleformat{\subsection}
  {\sffamily\bfseries\large\color{accent!85!black}\raggedright}{\thesubsection}{0.6em}{}
\titlespacing*{\subsection}{0pt}{1.2em plus 0.3em}{0.3em}
\titleformat{\paragraph}[runin]{\sffamily\bfseries}{}{0pt}{}[.]
\titlespacing*{\paragraph}{0pt}{0.6em}{0.5em}
\setcounter{secnumdepth}{2}
% a run-in heading that ends in its own punctuation (a question)
\newcommand{\qparagraph}[1]{\par\addvspace{0.6em}\noindent{\sffamily\bfseries #1}\enspace}
\setcounter{tocdepth}{1}

% ---------------------------------------------------------------- header
\pagestyle{fancy}
\fancyhf{}
\fancyhead[LE]{\small\sffamily\leftmark}
\fancyhead[RO]{\small\sffamily\rightmark}
\fancyfoot[C]{\small\sffamily\thepage}
\renewcommand{\headrulewidth}{0.3pt}
\renewcommand{\chaptermark}[1]{\markboth{\thechapter.\ #1}{}}
\renewcommand{\sectionmark}[1]{\markright{\thesection\ #1}}
\fancypagestyle{plain}{\fancyhf{}\fancyfoot[C]{\small\sffamily\thepage}\renewcommand{\headrulewidth}{0pt}}

% ---------------------------------------------------------------- inline code
% Grey background behind \texttt via lua-ul leaders: line spacing unchanged,
% line breaks still possible.
\usepackage{luacolor,lua-ul}
\NewCommandCopy{\origtexttt}{\texttt}
\DeclareRobustCommand{\texttt}[1]{\highLight[ttbg]{\origtexttt{#1}}}
% \code{...}: inline code; write \% for % (see companion.lua)
\directlua{companion = dofile("companion.lua")}
\DeclareRobustCommand{\code}[1]{\texttt{\directlua{companion.code("\luaescapestring{\detokenize{#1}}")}}}
\pdfstringdefDisableCommands{\def\code#1{#1}\def\texttt#1{#1}}

% ---------------------------------------------------------------- notation
\newcommand{\T}{^{\mkern-1.5mu\mathsf{T}}}
\newcommand{\inv}{^{-1}}
\newcommand{\invT}{^{-\mathsf{T}}}
\newcommand{\R}{\mathbb{R}}

% ---------------------------------------------------------------- links
% The pinned commit of ulm.lam that this text describes.
\newcommand{\commit}{41855425a8b9ea4ee279fd2ea4b31b45b53e1983}
\newcommand{\shortcommit}{4185542}
\newcommand{\repo}{https://github.com/michael-lehn/ulm-lam}
% \srcref{path}{first line}{last line}: link to the lines at the pinned commit
\newcommand{\srcref}[3]{\href{\repo/blob/\commit/#1\#L#2-L#3}{\code{#1:#2}}}
% \srcfile{path}: link to a file at the pinned commit
\newcommand{\srcfile}[1]{\href{\repo/blob/\commit/#1}{\code{#1}}}
% FG (Field Guide): \FG{1.7}{Refinement types} links to the named destination
% section.1.7 of modern-compiler-concepts.pdf (the two PDFs side by side);
% the printed text is the fallback.
\newcommand{\FGfile}{modern-compiler-concepts.pdf}
\newcommand{\FG}[2]{\href{\FGfile\#section.#1}{FG~§#1}~\emph{#2}}
\newcommand{\FGs}[1]{\href{\FGfile\#section.#1}{FG~§#1}}
\newcommand{\FGch}[2]{\href{\FGfile\#chapter.#1}{FG Chapter~#1}~\emph{#2}}
\newcommand{\FGchs}[1]{\href{\FGfile\#chapter.#1}{FG Chapter~#1}}

% ---------------------------------------------------------------- fact origins
\definecolor{odecl}{HTML}{C0661A}
\definecolor{over}{HTML}{1F5A96}
\definecolor{oinf}{HTML}{2E7D32}
\definecolor{oana}{HTML}{6A3D9A}
\definecolor{orun}{HTML}{B3261E}
\definecolor{onone}{HTML}{666666}
\newcommand{\otag}[2]{\highLight[#1!14]{\textsf{\small\color{#1!80!black}#2}}}
\newcommand{\Odecl}{\otag{odecl}{declared}}
\newcommand{\Over}{\otag{over}{verified}}
\newcommand{\Oinf}{\otag{oinf}{inferred}}
\newcommand{\Oana}{\otag{oana}{analysis}}
\newcommand{\Orun}{\otag{orun}{run time}}
\newcommand{\Onone}{\otag{onone}{not established}}

% ---------------------------------------------------------------- callouts
\newtcolorbox{theory}[1][]{enhanced,breakable,colback=theorybg,colframe=accent,boxrule=0pt,
  leftrule=2.5pt,arc=0pt,left=7pt,right=7pt,top=4pt,bottom=4pt,before skip=0.9em,after skip=0.9em,
  before upper={\textbf{\sffamily\color{accent}Theory#1.}\ }}
\newtcolorbox{limit}[1][]{enhanced,breakable,colback=limitbg,colframe=BurntOrange,boxrule=0pt,
  leftrule=2.5pt,arc=0pt,left=7pt,right=7pt,top=4pt,bottom=4pt,before skip=0.9em,after skip=0.9em,
  before upper={\textbf{\sffamily\color{BurntOrange!80!black}Not the same thing#1.}\ }}
\newtcolorbox{answer}{enhanced,breakable,colback=white,colframe=accent!60,boxrule=0.5pt,
  arc=1.5pt,left=7pt,right=7pt,top=4pt,bottom=4pt,before skip=0.6em,after skip=0.9em,
  before upper={\textbf{\sffamily\color{accent}Answer.}\ }}

% A step in the input / representation / rule / output / theory pattern.
\newenvironment{steps}
  {\par\begin{tcolorbox}[enhanced,colback=bg,colframe=bg,boxrule=0pt,arc=1.5pt,
      left=6pt,right=6pt,top=3pt,bottom=3pt,before skip=0.6em,after skip=0.9em]
   \begin{description}[font=\sffamily\bfseries\small,leftmargin=6.6em,labelwidth=6.2em,
      labelsep=0.4em,style=sameline,itemsep=2pt,topsep=0pt]}
  {\end{description}\end{tcolorbox}}

% ---------------------------------------------------------------- listings
\lstdefinelanguage{mlir}{
  alsoletter={.,_,\%,\#,!,@},
  morekeywords={func.func,return,func.call,scf.for,scf.if,scf.yield,arith.constant},
  morekeywords=[2]{ulm_lam.mul,ulm_lam.add,ulm_lam.scale,ulm_lam.transpose,
    ulm_lam.inverse,ulm_lam.identity,ulm_lam.assume,ulm_lamk.gemm,ulm_lamk.gemmtr,
    ulm_lamk.symm,ulm_lamk.trmm,ulm_lamk.trsm,ulm_lamk.syrk,ulm_lamk.gemv,ulm_lamk.trsv,
    ulm_lamk.potrf,ulm_lamk.getrf,ulm_lamk.getrs,ulm_lamk.gecopy,ulm_lamk.geaxpy,
    ulm_lamk.gescal,ulm_lamk.sycomplete,ulm_lamk.call,memref.copy,memref.view,memref.alloc,
    tensor.empty,bufferization.materialize_in_destination},
  morekeywords=[3]{tensor,memref,f64,i32,i64,f32,index},
  morecomment=[l]{//},
  morestring=[b]",
  sensitive=true}
\lstdefinelanguage{tablegen}{
  morekeywords={def,class,let,include,dag,string,list,code,bit,int},
  morekeywords=[2]{ins,outs},
  morecomment=[l]{//},
  morestring=[b]",
  sensitive=true}
\lstdefinelanguage{julia}{
  morekeywords={function,end,for,in,if,else,elseif,return,using,do,try,catch,let,const,struct,macro},
  morekeywords=[2]{@lam,@code_lam,Mat,Vec,Index,inv,assume,loop},
  morecomment=[l]{\#},
  morestring=[b]",
  sensitive=true}
\lstdefinestyle{base}{
  basicstyle=\ttfamily\footnotesize,
  keywordstyle=\color{kw}\bfseries,
  keywordstyle=[2]\color{kw},
  keywordstyle=[3]\color{ty},
  commentstyle=\color{cm}\itshape,
  stringstyle=\color{st},
  backgroundcolor=\color{bg},
  frame=none,
  xleftmargin=4pt, xrightmargin=4pt,
  framexleftmargin=4pt, framexrightmargin=4pt,
  aboveskip=6pt, belowskip=8pt,
  columns=fullflexible, keepspaces=true,
  showstringspaces=false, breaklines=true,
  breakindent=2em, postbreak=\mbox{\textcolor{cm}{$\hookrightarrow$}\space},
  upquote=true, tabsize=8, escapechar=`}
\lstdefinestyle{mlir}{style=base,language=mlir}
\lstdefinestyle{td}{style=base,language=tablegen}
\lstdefinestyle{cpp}{style=base,language=C++}
\lstdefinestyle{jl}{style=base,language=julia}
\lstdefinestyle{txt}{style=base,language={}}
\lstdefinestyle{sh}{style=base,language=bash}
\lstset{style=mlir}

% Tool output with Unicode (@code_lam: L1⁻ᵀ, ←): fvextra keeps the characters
% in order (listings moves modifier letters such as ᵀ before the identifier).
\usepackage{fvextra}
\newenvironment{textout}
  {\VerbatimEnvironment
   \begin{tcolorbox}[enhanced,breakable,colback=bg,colframe=bg,boxrule=0pt,arc=0pt,
      left=4pt,right=4pt,top=2pt,bottom=2pt,before skip=6pt,after skip=8pt]%
   \begin{Verbatim}[fontsize=\footnotesize,breaklines,breakanywhere,breaksymbolleft={}]}
  {\end{Verbatim}\end{tcolorbox}}

% ---------------------------------------------------------------- tables
\newcolumntype{L}{>{\raggedright\arraybackslash}X}
\newcolumntype{P}[1]{>{\raggedright\arraybackslash}p{#1}}
